\clearpage
\section{Phân phối có điều kiện: biết thêm một biến thì điều gì thay đổi?}
\label{sec:conditional}

\subsection{Phân phối biên và có điều kiện trả lời hai câu hỏi khác nhau}
Phân phối biên của chiều cao mô tả một người chọn ngẫu nhiên khi ta chưa
biết cân nặng. Phân phối chiều cao \emph{có điều kiện} tại cân nặng 70 kg
mô tả người trong nhóm có cân nặng đó. Nếu hai đặc điểm liên hệ, biết cân
nặng cung cấp thông tin về chiều cao và làm dự báo bớt bất định.

Ký hiệu $X_1\mid X_2=x_2$ không phải phép chia. Nó nói rằng ta xét phân
phối của $X_1$ sau khi biết giá trị của $X_2$. Vì biến liên tục có xác suất
đúng một điểm bằng 0, công thức có điều kiện được hiểu qua mật độ có điều
kiện, không tính bằng cách chia cho $\Prob(X_2=x_2)$.

\subsection{Công thức hai biến: dễ nhớ qua tương quan}
Với $(X_1,X_2)'$ chuẩn hai chiều, đặt
$\Var(X_1)=\sigma_1^2$, $\Var(X_2)=\sigma_2^2$ và tương quan $\rho$.
Khi $\sigma_2^2>0$,
\begin{align}
 \E(X_1\mid X_2=x_2)
 &=\mu_1+\rho\frac{\sigma_1}{\sigma_2}(x_2-\mu_2),
 \label{eq:cond2mean}\\
 \Var(X_1\mid X_2=x_2)&=\sigma_1^2(1-\rho^2).
 \label{eq:cond2var}
\end{align}
Toàn bộ phân phối có điều kiện là chuẩn với hai tham số này.
\nguon{Chương 4, §4.2, Result 4.6 và Example 4.7, tr.~160--162}

\hieu{Nếu $X_2$ cao hơn trung bình, và $\rho>0$, dự báo $X_1$ cũng dịch
lên. Độ dịch phụ thuộc tương quan và tỉ lệ đơn vị đo.
Thông tin về $X_2$ giải thích được phần $\rho^2$ của phương sai $X_1$;
phần còn lại là $1-\rho^2$.}

Nếu $\rho=0$, trung bình và phương sai không đổi: biết $X_2$ không giúp
dự báo $X_1$ trong mô hình chuẩn. Nếu $|\rho|$ gần 1, phương sai còn lại
nhỏ. Trường hợp $|\rho|=1$ là suy biến, không thuộc mô hình mật độ hai chiều
không suy biến đang xét.

\begin{example}[Một lần cập nhật dự báo]
Giả sử $\vmu=(10,20)'$, $\Sigma=\begin{pmatrix}4&3\\3&9\end{pmatrix}$.
Ta có $\rho=0.5$, $\sigma_1=2$, $\sigma_2=3$.
Khi quan sát $X_2=23$,
\[
 \E(X_1\mid X_2=23)=10+\frac39(23-20)=11,
\]
\[
 \Var(X_1\mid X_2=23)=4-\frac{3^2}{9}=3.
\]
Vậy $X_1\mid X_2=23\sim\N(11,3)$.
Trước khi biết $X_2$, $X_1\sim\N(10,4)$: dự báo đã dịch từ 10 đến 11,
độ lệch chuẩn giảm từ 2 xuống $\sqrt3$.
\end{example}

\begin{figure}[htbp]
\centering
\begin{tikzpicture}[x=.62cm,y=7cm,>=Stealth]
 \draw[->] (3,0)--(18,0) node[right]{$x_1$};
 \draw[->] (3,0)--(3,.28) node[above]{mật độ};
 \draw[domain=3:17,samples=90,smooth,dashed]
 plot (\x,{exp(-((\x-10)^2)/8)/sqrt(8*pi)});
 \draw[domain=3:17,samples=90,smooth]
 plot (\x,{exp(-((\x-11)^2)/6)/sqrt(6*pi)});
 \draw[dotted] (10,0)--(10,.2);
 \draw[dotted] (11,0)--(11,.231);
 \node[below] at (10,0) {10};
 \node[below] at (11,0) {11};
 \node[align=left,anchor=west] at (13,.27) {Nét đứt: biên $\N(10,4)$};
 \node[align=left,anchor=west] at (13,.18) {Nét liền: điều kiện $\N(11,3)$};
\end{tikzpicture}
\caption{Biết $X_2=23$ làm tâm dịch sang phải và phân phối $X_1$ hẹp lại
trong ví dụ. Diện tích dưới mỗi đường vẫn bằng 1.}
\end{figure}

\subsection{Công thức ma trận: dự báo một khối bằng một khối khác}
Phân hoạch vectơ và tham số thành
\[
 \mathbf X=\begin{pmatrix}\mathbf X_1\\\mathbf X_2\end{pmatrix},\quad
 \vmu=\begin{pmatrix}\vmu_1\\\vmu_2\end{pmatrix},\quad
 \Sigma=\begin{pmatrix}\Sigma_{11}&\Sigma_{12}\\\Sigma_{21}&\Sigma_{22}\end{pmatrix}.
\]
Khối $\mathbf X_1$ có $q_1$ thành phần; khối $\mathbf X_2$ có $q_2$ thành
phần. Chỉ số ``1'' ở đây là nhãn \emph{khối}, không nhất thiết chỉ biến
thứ nhất của dữ liệu.

\begin{result}
Nếu vectơ ghép chuẩn và $\Sigma_{22}$ khả nghịch, thì
\begin{equation}
 \mathbf X_1\mid\mathbf X_2=\mathbf x_2
 \sim\N_{q_1}\left(\vmu_{1\mid2},\Sigma_{11\mid2}\right),
 \label{eq:conditional}
\end{equation}
với
\begin{align}
 \vmu_{1\mid2}&=\vmu_1+\Sigma_{12}\Sigma_{22}^{-1}(\mathbf x_2-\vmu_2),
 \label{eq:conditionalmean}\\
 \Sigma_{11\mid2}&=\Sigma_{11}-\Sigma_{12}\Sigma_{22}^{-1}\Sigma_{21}.
 \label{eq:conditionalcov}
\end{align}
\end{result}
\nguon{Chương 4, §4.2, Result 4.6, tr.~160--161; dạng hồi quy (4-9), tr.~163}

Ma trận $B=\Sigma_{12}\Sigma_{22}^{-1}$ có kích thước $q_1\times q_2$
và đóng vai trò hệ số hồi quy. Nó biến độ lệch của khối đã biết thành phần
hiệu chỉnh cho khối cần dự báo.
Khối trừ đi trong \eqref{eq:conditionalcov} là phần phương sai được giải
thích bởi thông tin đó; $\Sigma_{11\mid2}$ còn được gọi là bổ Schur.

Đặc biệt, hiệp phương sai có điều kiện không phụ thuộc \emph{giá trị}
$\mathbf x_2$. Thay $\mathbf x_2$ làm tâm dịch chuyển nhưng không làm
độ phân tán có điều kiện thay đổi. Tính chất này là của mô hình chuẩn;
không suy ra như vậy cho mọi phân phối.

\subsection{Quy trình tính: ghi đúng khối trước khi nhân}
\begin{example}[Dự báo một biến từ hai biến]
Với \eqref{eq:running3d}, tìm $X_1\mid(X_2=x_2,X_3=x_3)$.
\buoc{1}{Đặt khối mục tiêu và khối đã biết}
Đặt $U=X_1$, $\mathbf V=(X_2,X_3)'$.
\buoc{2}{Trích tham số đúng thứ tự}
\[
 \mu_U=1,\quad\vmu_V=\begin{pmatrix}-1\\2\end{pmatrix},\quad
 \Sigma_{UU}=4,\quad\Sigma_{UV}=(0,-1),\quad
 \Sigma_{VV}=\begin{pmatrix}5&0\\0&2\end{pmatrix}.
\]
\buoc{3}{Tính hệ số cập nhật}
\[
 B=(0,-1)\begin{pmatrix}1/5&0\\0&1/2\end{pmatrix}=(0,-1/2).
\]
\buoc{4}{Tính hai tham số}
\[
 \E(U\mid\mathbf V)=1+(0,-1/2)
 \begin{pmatrix}x_2+1\\x_3-2\end{pmatrix}=2-\frac{x_3}{2},
\]
\[
 \Var(U\mid\mathbf V)=4-(0,-1/2)(0,-1)'=\frac72.
\]
Vậy $X_1\mid(X_2=x_2,X_3=x_3)\sim\N(2-x_3/2,7/2)$.
Giá trị $x_2$ biến mất vì $X_2$ độc lập với cả $(X_1,X_3)'$.
\end{example}

Nếu chỉ biết $X_3=x_3$, ta vẫn nhận được đúng phân phối trên. Khi
$x_3=4$, trung bình có điều kiện bằng 0. Không có mâu thuẫn với trung bình
biên bằng 1: đó là trung bình trên tất cả các giá trị của $X_3$, còn ở đây
ta đã giới hạn vào một nhóm có $X_3=4$.

\subsection{Phần dư hồi quy và tính độc lập}
Đặt phần dư đã trừ trung bình
\[
 \mathbf R=\mathbf X_1-\vmu_1-B(\mathbf X_2-\vmu_2),
 \qquad B=\Sigma_{12}\Sigma_{22}^{-1}.
\]
Ta có $\E(\mathbf R)=0$ và
\[
 \Cov(\mathbf R,\mathbf X_2)=\Sigma_{12}-B\Sigma_{22}=0.
\]
Vì $(\mathbf R,\mathbf X_2)$ chuẩn đồng thời, $\mathbf R\perp\mathbf X_2$.
Do đó khi biết $\mathbf X_2$, phần dư vẫn có cùng phân phối; điều kiện hóa
chỉ biến $B(\mathbf X_2-\vmu_2)$ thành một độ dịch tâm xác định.
Đây là cách suy ra công thức có điều kiện mà không cần tính tích phân.

\luuy{Ở đây các hệ số $B$ tính từ tham số tổng thể đã biết. Đây là tính
chất của phân phối, chưa phải bài toán ước lượng hệ số hồi quy từ mẫu.
Không tự thay bậc tự do của hồi quy vào công thức này.}

\kiemtra{Có thể lấy ma trận hiệp phương sai biên $\Sigma_{11}$ làm
hiệp phương sai có điều kiện không? Chỉ khi phần hiệu chỉnh bằng 0,
chẳng hạn hai khối độc lập.}
